\documentclass[opre]{informs3}
\OneAndAHalfSpacedXI

\usepackage[utf8]{inputenc}
\usepackage{amsmath,amssymb,mathrsfs}
\usepackage{bm,bbm,booktabs,float,graphicx}
\usepackage{natbib}
\bibpunct[, ]{(}{)}{,}{a}{}{,}%
\def\bibfont{\small}%
\def\bibsep{\smallskipamount}%
\def\bibhang{24pt}%
\def\newblock{\ }%
\def\BIBand{and}%
\usepackage[hyperfootnotes=false,colorlinks=true,linkcolor=black,urlcolor=black,citecolor=blue]{hyperref}
\usepackage{cleveref}

\TheoremsNumberedThrough
\EquationsNumberedThrough
\allowdisplaybreaks

\newcommand{\E}{\mathbb{E}}
\newcommand{\Pb}{\mathbb{P}}
\newcommand{\Ind}[1]{\mathbbm{1}\{#1\}}
\newcommand{\Xscr}{\mathcal{X}}

\graphicspath{{../outputs_gamma_0999/plots/}{../outputs/plots/}}

\begin{document}

\RUNAUTHOR{Natalia Trigo}
\RUNTITLE{Discounted Dynamic Reputation Investment}

\TITLE{Discounted Dynamic Reputation Investment with Hidden Product Choice}

\ARTICLEAUTHORS{
\AUTHOR{Natalia Trigo}
\AFF{Anderson UCLA, \EMAIL{natalia.trigo.phd@anderson.ucla.edu}}
}

\ABSTRACT{
We study a discounted reputation problem in which a seller privately chooses
between a low-cost product and a high-quality product while a user learns only
from realized binary outcomes. The user treats the seller as a single Bernoulli
arm and applies Thompson sampling to decide whether to engage with the seller or
with a known outside option. The seller chooses hidden products to maximize
expected discounted profit. We formulate the infinite-horizon Bellman equation
and solve a long-horizon temporal approximation by exact backward induction over
all reachable posterior states. The implementation imposes no artificial
boundary on the number of observations. The optimal product choice is governed
by the discounted continuation value of a success relative to a failure. We
also derive a deterministic discounted fluid approximation that preserves this
threshold structure and identifies a long-run quality-mix benchmark. Numerical
experiments with a patient seller illustrate how reputation investment varies
with the quality of the outside option.
}

\KEYWORDS{dynamic reputation; hidden actions; discounted dynamic programming; fluid approximation; Thompson sampling}

\maketitle

\section{Model Formulation}

Time is discrete and indexed by $t=1,2,\ldots$. The environment consists of a
single user and two sellers, denoted by $A$ and $B$. In each period, the user
selects a seller $y_t\in\{A,B\}$ and receives binary utility $u_t\in\{0,1\}$.

Seller B is a transparent outside option. If the user selects Seller B, utility
is Bernoulli with known success probability $p_0\in(0,1)$:
\begin{equation}
    \Pb(u_t=1\mid y_t=B)=p_0.
    \label{eq:seller-b-success}
\end{equation}

Seller A has two products, $\Xscr=\{1,2\}$. Conditional on being selected,
Seller A privately assigns product $x_t\in\Xscr$. Product $i$ generates a
successful outcome with probability $p_i$ and costs Seller A $c_i$:
\begin{equation}
    \Pb(u_t=1\mid y_t=A,x_t=i)=p_i.
    \label{eq:seller-a-success}
\end{equation}
We assume
\begin{equation}
    0<p_1<p_2\leq 1,\qquad 0\leq c_1<c_2,
    \label{eq:product-ordering}
\end{equation}
so product 2 is more likely to satisfy the user but is also more costly.
Whenever Seller A is chosen, it receives fixed revenue $R>0$. Its period profit
is
\begin{equation}
    \pi_t^A=
    \begin{cases}
        R-c_{x_t}, & y_t=A,\\
        0, & y_t=B.
    \end{cases}
    \label{eq:period-profit}
\end{equation}

\section{User Learning and Demand}

The user does not observe Seller A's product choice. Product quality affects the
distribution of realized utility, but the user sees only the selected seller
and the binary outcome.

\begin{assumption}[User's Misspecified Belief]
After selecting Seller A, the user observes $u_t$ but not $x_t$. The user
therefore treats Seller A as a single product with an unknown Bernoulli success
probability $\theta\in[0,1]$.
\end{assumption}

Let $S_t$ and $F_t$ be the cumulative successes and failures generated by
Seller A before period $t$:
\begin{align}
    S_t &= \sum_{\tau=1}^{t-1}\Ind{y_\tau=A,u_\tau=1},\\
    F_t &= \sum_{\tau=1}^{t-1}\Ind{y_\tau=A,u_\tau=0}.
\end{align}
Starting from the uniform prior
$\theta\sim\operatorname{Beta}(1,1)$, the posterior is
\begin{equation}
    \theta\mid S_t,F_t
    \sim
    \operatorname{Beta}(1+S_t,1+F_t).
    \label{eq:posterior}
\end{equation}
The corresponding posterior mean is
\begin{equation}
    m(S_t,F_t)=\frac{S_t+1}{S_t+F_t+2}.
    \label{eq:posterior-mean}
\end{equation}

The user follows Thompson sampling. At the beginning of period $t$, the user
draws
\begin{equation}
    \widetilde{\theta}_t
    \sim
    \operatorname{Beta}(1+S_t,1+F_t)
\end{equation}
and selects Seller A if and only if
$\widetilde{\theta}_t\geq p_0$. Seller A observes the user's choices and the
realized outcomes, and can therefore track $(S_t,F_t)$. At state $(S,F)$, the
probability that Seller A is selected is
\begin{equation}
\begin{aligned}
    D(S,F)
    &\triangleq
    \Pb(\widetilde{\theta}_t\geq p_0\mid S_t=S,F_t=F)\\
    &=
    \int_{p_0}^{1}
    \frac{\theta^S(1-\theta)^F}{B(1+S,1+F)}\,d\theta,
\end{aligned}
\label{eq:demand-probability}
\end{equation}
where $B(\cdot,\cdot)$ is the Beta function. Demand determines how often
Seller A earns revenue and receives a new reputation signal.

\section{Discounted Seller Problem}

Let $\gamma\in(0,1)$ be Seller A's discount factor. Seller A chooses a policy
that maps the user's posterior state to a hidden product and maximizes
\begin{equation}
    \E\left[
        \sum_{t=1}^{\infty}
        \gamma^{t-1}
        \Ind{y_t=A}(R-c_{x_t})
    \right].
    \label{eq:discounted-objective}
\end{equation}
Let $V(S,F)$ be the optimal infinite-horizon value. If the user chooses Seller
B, Seller A earns no current profit, the state does not change, and calendar
time still advances. The stationary Bellman equation is
\begin{equation}
\begin{aligned}
    V(S,F)
    ={}&
    \gamma[1-D(S,F)]V(S,F)\\
    &+
    D(S,F)
    \max_{x\in\{1,2\}}
    \left\{
        R-c_x
        +\gamma\left[
            p_xV(S+1,F)+(1-p_x)V(S,F+1)
        \right]
    \right\}.
\end{aligned}
\label{eq:stationary-bellman}
\end{equation}
This equation characterizes the discounted model. The computation below does
not substitute a finite-horizon policy for the discounted policy. Instead, it
approximates the solution to \cref{eq:stationary-bellman} by retaining a long
discounted tail and then checks whether the early-period policy is stable as
the computational horizon increases.

\subsection{Exact Long-Horizon Backward Induction}

For a computational horizon $T$, let $V_t^T(S,F)$ denote the discounted value
from periods $t,\ldots,T$, with
\begin{equation}
    V_{T+1}^T(S,F)=0.
    \label{eq:terminal-condition}
\end{equation}
The exact recursion is
\begin{equation}
\begin{aligned}
    V_t^T(S,F)
    ={}&
    \gamma[1-D(S,F)]V_{t+1}^T(S,F)\\
    &+
    D(S,F)
    \max_{x\in\{1,2\}}
    \big\{
        R-c_x
        +\gamma[
            p_xV_{t+1}^T(S+1,F)
            +(1-p_x)V_{t+1}^T(S,F+1)
        ]
    \big\}.
\end{aligned}
\label{eq:discounted-bellman}
\end{equation}
At period $t$, every reachable state satisfies $S+F\leq t-1$. Backward
induction therefore covers the complete triangular state space for the stated
horizon. Only adjacent value layers are stored in memory. In particular, the
algorithm does not truncate at a fixed value of $S+F$, make a count boundary
absorbing, or project simulated states back onto a smaller state space.

Conditional on the user choosing Seller A, define
\begin{equation}
    Q_{t,x}^T(S,F)
    \triangleq
    R-c_x
    +\gamma[
        p_xV_{t+1}^T(S+1,F)
        +(1-p_x)V_{t+1}^T(S,F+1)
    ].
    \label{eq:action-value}
\end{equation}
Let the raw continuation difference be
\begin{equation}
    G_t^T(S,F)
    \triangleq
    V_{t+1}^T(S+1,F)-V_{t+1}^T(S,F+1)
    \label{eq:raw-gap}
\end{equation}
and define the discounted continuation value gap
\begin{equation}
    M_t^T(S,F)
    \triangleq
    \gamma G_t^T(S,F).
    \label{eq:discounted-gap}
\end{equation}
Writing $\Delta p=p_2-p_1$ and $\Delta c=c_2-c_1$, the action-value difference
is
\begin{equation}
    Q_{t,2}^T(S,F)-Q_{t,1}^T(S,F)
    =
    -\Delta c+\Delta p\,M_t^T(S,F).
    \label{eq:q-gap}
\end{equation}
Consequently, product 2 is optimal if and only if
\begin{equation}
    M_t^T(S,F)\geq\frac{\Delta c}{\Delta p}.
    \label{eq:discounted-threshold}
\end{equation}
If the raw gap $G_t^T$ is reported instead, its threshold is
$\Delta c/(\gamma\Delta p)$. The distinction matters when interpreting value
gap plots; the reported object $M_t^T$ already contains the discount factor.

\section{Discounted Fluid Approximation}

The exact dynamic program retains the full distribution of the posterior count
process. A deterministic fluid approximation replaces random increments by
their conditional expected drifts. This provides an interpretable long-run
benchmark and avoids enumerating a growing triangular state space. It does not
alter the information structure: the user still observes outcomes but never
observes Seller A's product.

\subsection{State Dynamics and Objective}

Let $\tau\geq0$ denote continuous calendar time and let $s(\tau)$ and $f(\tau)$
be fluid success and failure counts. Extend demand to nonnegative real counts:
\begin{equation}
    \mathcal{D}(s,f)
    \triangleq
    \int_{p_0}^{1}
    \frac{\theta^s(1-\theta)^f}{B(1+s,1+f)}\,d\theta.
    \label{eq:fluid-demand}
\end{equation}
Let $a(\tau)\in[0,1]$ be the relaxed intensity of product 2 and define
\begin{align}
    q(a)&=p_1+a\Delta p,\\
    c(a)&=c_1+a\Delta c.
\end{align}
The expected count dynamics are
\begin{align}
    \dot{s}(\tau)
    &=
    \mathcal{D}(s(\tau),f(\tau))q(a(\tau)),
    \label{eq:fluid-success}\\
    \dot{f}(\tau)
    &=
    \mathcal{D}(s(\tau),f(\tau))[1-q(a(\tau))].
    \label{eq:fluid-failure}
\end{align}
Set $r=-\log\gamma$, so the continuous-time discount factor over $\tau$ periods
is $e^{-r\tau}=\gamma^\tau$. For horizon $H$, the discounted fluid objective is
\begin{equation}
    J_H^{\mathrm{fl}}(a)
    =
    \int_0^H
    e^{-r\tau}
    \mathcal{D}(s(\tau),f(\tau))
    [R-c(a(\tau))]\,d\tau.
    \label{eq:fluid-objective}
\end{equation}
This is a mean-drift closure, not an exact reformulation: in general,
$\E[D(S_t,F_t)]\neq D(\E[S_t],\E[F_t])$.

\subsection{Fluid Optimality Condition}

Let $\mathcal{V}(\tau,s,f)$ be the optimal fluid value from $\tau$ to $H$,
discounted relative to time $\tau$, with $\mathcal{V}(H,s,f)=0$. Its
Hamilton--Jacobi--Bellman equation is
\begin{equation}
\begin{aligned}
    -\partial_\tau\mathcal{V}+r\mathcal{V}
    =
    \mathcal{D}(s,f)
    \max_{a\in[0,1]}
    \big\{
        R-c(a)
        +q(a)\partial_s\mathcal{V}
        +[1-q(a)]\partial_f\mathcal{V}
    \big\}.
\end{aligned}
\label{eq:fluid-hjb}
\end{equation}
The coefficient multiplying $a$ inside the Hamiltonian is
\begin{equation}
    -\Delta c
    +\Delta p
    (\partial_s\mathcal{V}-\partial_f\mathcal{V}).
    \label{eq:fluid-control-coefficient}
\end{equation}
Thus the relaxed solution is bang-bang except at indifference. Product 2 is
optimal whenever
\begin{equation}
    \partial_s\mathcal{V}-\partial_f\mathcal{V}
    \geq
    \frac{\Delta c}{\Delta p}.
    \label{eq:fluid-threshold}
\end{equation}
The fluid model therefore preserves the discrete model's central tradeoff:
the costly product is used when the shadow value of replacing a failure by a
success is sufficiently large. Discounting changes the value function and its
derivatives through $r\mathcal{V}$, even though it does not appear separately
in the threshold once the discounted shadow value has been computed.

\subsection{Effective Observations and Long-Run Quality Mix}

Define effective observations and the fluid posterior mean by
\begin{equation}
    n(\tau)=s(\tau)+f(\tau),
    \qquad
    m(\tau)=\frac{s(\tau)+1}{n(\tau)+2}.
    \label{eq:fluid-transform}
\end{equation}
Let $\widehat{\mathcal{D}}(n,m)$ denote \cref{eq:fluid-demand} written in this
transformed state. Then
\begin{align}
    \dot{n}(\tau)
    &=
    \widehat{\mathcal{D}}(n(\tau),m(\tau)),
    \label{eq:fluid-count}\\
    \dot{m}(\tau)
    &=
    \frac{\widehat{\mathcal{D}}(n(\tau),m(\tau))}{n(\tau)+2}
    [q(a(\tau))-m(\tau)].
    \label{eq:fluid-mean}
\end{align}
Demand determines the rate at which the seller collects observations, while
the delivered quality determines the direction of belief updating. The factor
$1/(n+2)$ captures belief rigidity as evidence accumulates.

If a constant control $a$ is used and demand remains positive, dividing
\cref{eq:fluid-mean} by \cref{eq:fluid-count} gives
\begin{equation}
    \frac{dm}{dn}=\frac{q(a)-m}{n+2}.
\end{equation}
Starting from the uniform prior, the posterior mean converges to $q(a)$. When
$p_1<p_0\leq p_2$, the product-2 intensity that places limiting quality exactly
at the outside option is
\begin{equation}
    a_{\mathrm{bdry}}(p_0)
    =
    \frac{p_0-p_1}{p_2-p_1}.
    \label{eq:boundary-mix}
\end{equation}
For $p_0\leq p_1$, product 1 alone clears this benchmark. For $p_0>p_2$, no
feasible mix can sustain a limiting quality at least as high as the outside
option. The boundary mix is a sustainability benchmark, not the solution to
the discounted control problem: optimal investment also reflects current
costs, discounting, posterior uncertainty, and the option value of future
demand.

\section{Numerical Study}

\subsection{Calibration and Computation}

The numerical study uses the calibration in \cref{tab:calibration}. Product 2
raises the success probability by $\Delta p=0.45$ and costs an additional
$\Delta c=0.60$. Hence the discounted continuation-gap threshold in
\cref{eq:discounted-threshold} is
\begin{equation}
    \frac{\Delta c}{\Delta p}=\frac{4}{3}.
\end{equation}

\begin{table}[H]
\TABLE
{Numerical calibration.\label{tab:calibration}}
{\begin{tabular}{lccccc}
\toprule
Parameter & $p_1$ & $p_2$ & $c_1$ & $c_2$ & $R$\\
\midrule
Value & 0.35 & 0.80 & 0.05 & 0.65 & 1.00\\
\bottomrule
\end{tabular}}
{}
\end{table}

The main patient-seller experiment sets $\gamma=0.999$, solves
\cref{eq:discounted-bellman} through $T=2000$, and simulates the first 250
periods for 400 independent replications at each
$p_0\in\{0.1,0.3,0.5,0.7,0.9\}$. The final-window statistics use periods
201--250.

For $\gamma=0.98$, the default horizon $T=700$ satisfies
$\gamma^T=7.22\times10^{-7}$, so the omitted discounted tail is negligible
under a $10^{-6}$ criterion. For $\gamma=0.999$, however,
$\gamma^{2000}=0.1352$; a literal $\gamma^T<10^{-6}$ requirement would need
$T=13809$. Thus the $T=2000$ exercise should be interpreted as a
large-horizon policy experiment, not as a strict value-tail approximation.
The appropriate diagnostic is stability of the early policy under further
horizon increases.

\subsection{Simulation Results}

\Cref{fig:simulation} reports market share, product use, posterior beliefs, and
profit over the simulation window. The late-window results are summarized in
\cref{tab:late-results}.

\begin{figure}[H]
\centering
\includegraphics[width=\textwidth]{discounted_simulation_by_period.png}
\caption{Simulation outcomes under the exact discounted policy with
$\gamma=0.999$ and $T=2000$. Lines aggregate 400 independent replications for
each outside-option quality.}
\label{fig:simulation}
\end{figure}

\begin{table}[H]
\TABLE
{Late-window simulation results, periods 201--250.\label{tab:late-results}}
{\begin{tabular}{ccccc}
\toprule
$p_0$ &
Fluid boundary $a_{\mathrm{bdry}}$ &
Seller A share &
Product-2 rate &
Posterior mean\\
\midrule
0.10 & 0     & 1.000 & 0.000 & 0.354\\
0.30 & 0     & 0.992 & 0.009 & 0.387\\
0.50 & 0.333 & 0.994 & 0.420 & 0.588\\
0.70 & 0.778 & 0.935 & 0.883 & 0.763\\
0.90 & infeasible & 0.051 & 0.998 & 0.686\\
\bottomrule
\end{tabular}}
{\footnotesize The product-2 rate is conditional on Seller A being selected.
The fluid boundary is the minimum quality-mix benchmark in
\cref{eq:boundary-mix}; it is not an optimality prediction.}
\end{table}

When $p_0\in\{0.1,0.3\}$, the low-cost product already has quality at least as
high as the outside option, and reputation investment vanishes in the late
window. At $p_0=0.5$, Seller A is retained in 99.4\% of interactions and uses
product 2 in 42.0\% of its selected periods. At $p_0=0.7$, stronger quality
investment keeps Seller A's market share near 93.5\%. The optimal product-2
rates exceed the fluid boundary at both interior cutoffs, consistent with the
fact that equality only centers the limiting posterior at the user's cutoff.
At $p_0=0.9>p_2$, even product 2 cannot match the outside option. Seller A uses
it almost whenever selected, but its late market share falls to 5.1\%.

\Cref{fig:fluid-benchmark} compares the simulated product-2 rates with the fluid
quality-maintenance benchmark. The comparison is qualitative: stochastic
posterior paths, endogenous selection, and discounted option value are absent
from the deterministic boundary calculation.

\begin{figure}[H]
\centering
\includegraphics[width=\textwidth]{discounted_product2_paths_with_quality_benchmark.png}
\caption{Product-2 use under the exact discounted policy and the clipped fluid
quality-mix benchmark. The benchmark indicates the product mix required to
match the outside-option quality when such a mix is feasible.}
\label{fig:fluid-benchmark}
\end{figure}

Finally, \cref{fig:policy-heatmaps} displays early-period policy slices. Yellow
states use product 1 and teal states use product 2. The switching region is not
a rule based only on the posterior mean: the total number of observations and
the remaining discounted continuation value also affect whether an additional
success is worth its cost.

\begin{figure}[H]
\centering
\includegraphics[width=0.85\textwidth]{best_response_policy_heatmaps.png}
\caption{Exact discounted product policy over reachable posterior states. The
policy uses product 2 precisely where the discounted continuation value gap
exceeds $4/3$.}
\label{fig:policy-heatmaps}
\end{figure}

\section{Discussion}

The model separates three objects that are easy to conflate. Demand
$D(S,F)$ determines how frequently Seller A earns revenue and updates the
user's belief. The posterior state determines how responsive reputation is to
another outcome. The discounted gap $M_t^T(S,F)$ measures the economic value of
turning the next failure state into the next success state. Product 2 is chosen
only when that last object is large enough to cover its incremental cost.

The fluid approximation preserves this mechanism through a shadow-value
condition and adds a transparent benchmark for the long-run product mix.
However, it suppresses path dispersion and is least reliable near
$m=p_0$, where small posterior fluctuations sharply change Thompson-sampling
demand. It should therefore complement, rather than replace, exact backward
induction and stochastic simulation.

\section{Conclusion}

Reputation investment varies nonmonotonically with the user's outside option.
When the outside option is weak, the cheap product is sufficient. At
intermediate cutoffs, the seller uses the costly product to keep its posterior
reputation above the user's threshold. When the outside option dominates even
the high-quality product, quality investment cannot prevent demand loss. The
exact discounted recursion and the fluid benchmark expose these forces from
complementary perspectives.

\end{document}
